% Seção: Introdução  (orçamento: 1 página)
%
% Curta e sem detalhe: o como fica na Metodologia e os números nos Resultados. Sem hipóteses formais (ALERTAS,
% item 2) e sem parágrafo de organização do texto.

\section{Introdução}\label{sec:introdução}

Os dados de eventos de futebol registram cada ação de uma partida, com o jogador, o momento e a posição no
campo, e há conjuntos desse tipo abertos ao público~\cite{pappalardo2019public}. Muitas perguntas de análise
sobre esses dados exigem ligar eventos entre si: quem mais passou para quem, em que ordem a bola circulou em uma
jogada, o que mudou depois de uma substituição.

Um LLM permite fazer essas perguntas em linguagem natural, mas os dados da partida precisam chegar ao seu
contexto. Isso pode ser feito de várias formas: com a tabela inteira no prompt, com as linhas mais parecidas com
a pergunta, com estatísticas resumidas, com uma consulta escrita pelo próprio modelo ou com ferramentas que já
trazem a consulta pronta. Em todas, o modelo não é treinado de novo; muda o que entra no contexto e quem junta
os eventos, o modelo ou o código. No esporte, os trabalhos publicados avaliam uma dessas formas de cada vez
(Seção~\ref{sec:trabalhos}).

Este trabalho compara \nArms{} abordagens sobre a final da Copa do Mundo de 2022, com \nQuestions{} perguntas
cujo gabarito foi calculado dos dados brutos, em \nModels{} modelos de tamanhos diferentes. As perguntas de
pesquisa são qual abordagem acerta mais, como o acerto muda de um modelo para outro e quanto cada abordagem
custa em tokens e tempo. \armGraphTools{} teve acerto alto em todos os modelos. \armTextToCypher{} e
\armEventsInPrompt{} só se aproximaram dela nos modelos maiores, com mais tokens de saída e mais tempo por
pergunta.
